\chapter*{Introduction générale}
\addcontentsline{toc}{chapter}{Introduction générale}
\markboth{Introduction générale}{Introduction générale}

% ---------------------------------------------------------------
\section*{Contexte général}
\addcontentsline{toc}{section}{Contexte général}
% ---------------------------------------------------------------

Dans un environnement économique en constante mutation, les organisations sont soumises à des pressions croissantes pour garantir la conformité réglementaire, maîtriser leurs risques opérationnels et saisir les opportunités d'amélioration continue. Les référentiels normatifs internationaux, tels que l'ISO 9001 (qualité), l'ISO 45001 (santé-sécurité au travail) et l'ISO 14001 (environnement), imposent désormais une gestion formalisée et proactive des risques et des opportunités (R\&O) comme condition indispensable à la certification et au maintien de la performance.

La digitalisation des systèmes de management QHSE (Qualité, Hygiène, Sécurité, Environnement) a permis de centraliser ces données dans des plateformes spécialisées, à l'image de QALITAS développée par TIM Group Tunisie. Ces outils offrent aux responsables qualité des tableaux de bord riches en indicateurs de performance (KPI), des registres de non-conformités (NC), des cartographies de risques et des plans d'actions. Cependant, malgré cette richesse informationnelle, le pilotage des R\&O reste largement manuel, fragmenté et réactif.

Parallèlement, l'émergence des grands modèles de langage (LLM — Large Language Models) et des architectures multi-agents a ouvert de nouvelles perspectives pour l'automatisation de tâches à forte valeur cognitive : analyse contextuelle, raisonnement structuré, synthèse décisionnelle et génération d'actions priorisées. Ces technologies, récemment encadrées par des frameworks d'orchestration tels que LangGraph, permettent aujourd'hui de concevoir des systèmes intelligents capables d'opérer de façon autonome et coordonnée sur des données métier complexes.

% ---------------------------------------------------------------
\section*{Problématique}
\addcontentsline{toc}{section}{Problématique}
% ---------------------------------------------------------------

Malgré la richesse des données disponibles sur des plateformes comme QALITAS, les responsables qualité sont confrontés à plusieurs limites structurelles dans la gestion continue de leurs R\&O :

\begin{itemize}
    \item \textbf{Monitoring discontinu :} La réévaluation des risques est ponctuelle et dépend de l'initiative humaine, sans mécanisme d'alerte automatique sur les dépassements de seuils critiques.

    \item \textbf{Sous-exploitation des données :} Les KPI hors cible, les non-conformités récurrentes et les tendances négatives disponibles dans les tableaux de bord ne sont pas systématiquement corrélés aux entrées risques correspondantes.

    \item \textbf{Réévaluation tardive :} Le score de criticité RPN (Risk Priority Number), calculé à partir de la probabilité, de l'impact et de la détectabilité, n'est pas automatiquement mis à jour lorsque les conditions métier évoluent.

    \item \textbf{Opportunités non capitalisées :} Les processus performants ou les tendances favorables ne sont pas exploités de façon systématique pour générer des opportunités d'amélioration conformes aux exigences ISO.

    \item \textbf{Actions sans contexte :} Les plans d'actions correctifs ou préventifs générés manuellement manquent souvent de priorisation, d'argumentation et de lien explicite avec les risques identifiés.
\end{itemize}

Face à ces constats, la question centrale qui guide ce travail est la suivante :

\begin{center}
\textit{Comment concevoir un système multi-agents intelligent capable d'extraire, d'analyser et d'exploiter de manière autonome les données des tableaux de bord qualité pour assurer le monitoring continu, la réévaluation dynamique et le pilotage proactif des risques et des opportunités dans un système de management QHSE ?}
\end{center}

% ---------------------------------------------------------------
\section*{Solution proposée}
\addcontentsline{toc}{section}{Solution proposée}
% ---------------------------------------------------------------

Pour répondre à cette problématique, nous proposons la conception et le développement d'une \textbf{plateforme d'agents IA pour la gestion des risques et opportunités des systèmes de management QALITAS QSE}. Cette plateforme repose sur une architecture multi-agents composée de quatre agents spécialisés, coordonnés par un orchestrateur central implémenté avec le framework LangGraph.

Chaque agent est responsable d'une étape précise du cycle de pilotage des R\&O, depuis l'extraction et la structuration des données brutes jusqu'à l'injection automatique des actions correctives dans QALITAS via son API REST. Les agents communiquent à travers un état partagé (\textit{PipelineState}) et peuvent déclencher des boucles de réévaluation automatiques lorsque le niveau de criticité d'un risque évolue significativement.

La plateforme s'appuie sur plusieurs sources de données hétérogènes : l'API REST de QALITAS, les exports Excel des cartographies de risques, les fichiers CSV d'indicateurs de performance (KPI) et les rapports PDF des tableaux de bord qualité. Elle intègre des modèles de langage (LLM) pour l'enrichissement sémantique, la génération d'arguments contextualisés et la production de synthèses lisibles par les responsables qualité.

% ---------------------------------------------------------------
\section*{Plan du rapport}
\addcontentsline{toc}{section}{Plan du rapport}
% ---------------------------------------------------------------

Le présent rapport est structuré en quatre chapitres complémentaires, qui décrivent l'ensemble du cycle de vie du projet, de la présentation de l'environnement de travail jusqu'à la validation des résultats sur données réelles.

\vspace{0.8em}
\noindent\textbf{Chapitre 1 — Présentation de l'entreprise et cahier des charges}

Ce chapitre introduit le contexte général du stage réalisé au sein de TIM Group Tunisie. Nous y présentons l'entreprise d'accueil (historique, atouts, chiffres clés, services et références clients), le produit QALITAS sur lequel porte le projet, ainsi que le contexte qui a motivé ce travail. Nous détaillons ensuite le cahier des charges fonctionnel et technique de la plateforme, en décrivant les quatre agents du système, leurs rôles respectifs et les exigences transverses (robustesse, traçabilité, conformité API).

\vspace{0.8em}
\noindent\textbf{Chapitre 2 — État de l'art et étude de l'existant}

Ce chapitre présente une revue des connaissances nécessaires à la compréhension et à la justification des choix technologiques effectués. Nous abordons successivement : les fondements de la gestion des risques et opportunités dans les systèmes de management ISO, les approches existantes de monitoring qualité, les architectures multi-agents et leurs applications en milieu industriel, ainsi que les grands modèles de langage (LLM) et les frameworks d'orchestration agentique (LangGraph, LangChain). Nous étudions également l'existant QALITAS : structure de l'API REST, modèles de données, endpoints disponibles et contraintes d'intégration.

\vspace{0.8em}
\noindent\textbf{Chapitre 3 — Analyse et conception du système multi-agents}

Ce chapitre est consacré à la conception détaillée de la plateforme. Nous y présentons l'architecture globale du système (graphe LangGraph, flux de données, état partagé \textit{PipelineState}), les diagrammes UML de cas d'utilisation, de séquence et de classes. Nous décrivons la conception de chacun des quatre agents : Agent 1 (identification et caractérisation des R\&O), Agent 2 (analyse et réévaluation du score RPN), Agent 3 (génération des plans de traitement et détection des opportunités) et Agent 4 (monitoring, alertes et injection dans QALITAS). Nous détaillons également les mécanismes de boucle de réévaluation automatique et de gestion du cache d'injection.

\vspace{0.8em}
\noindent\textbf{Chapitre 4 — Réalisation et résultats}

Ce dernier chapitre présente l'implémentation concrète de la plateforme et les résultats obtenus sur des données réelles issues de QALITAS. Nous décrivons l'environnement technologique (Python, LangGraph, Ollama, API QALITAS), les choix d'implémentation de chaque agent, et les tests de validation réalisés. Nous analysons les résultats de l'injection automatique des actions dans QALITAS (risques couverts, actions générées, score RPN réévalué), et nous évaluons la qualité et la pertinence des sorties produites par les LLM. Une discussion critique des limites et des perspectives d'amélioration clôt ce chapitre.

\vspace{0.8em}
\noindent Le rapport se termine par une \textbf{conclusion générale} qui synthétise les apports du projet, évalue l'atteinte des objectifs fixés dans le cahier des charges, et trace des perspectives concrètes pour l'évolution future de la plateforme : intégration temps-réel, apprentissage continu, extension à d'autres modules QALITAS et déploiement en environnement de production.
